%% manuscript_v0_submitted.tex
%% LaTeX transcription of the manuscript as submitted to IMMI
%% (original_submission/CrabNet_Hyperparameters_IMMI_submission.docx),
%% used as the latexdiff baseline for the tracked-changes PDF.
%% Text, headings, figures and captions follow the submitted Word file.
%% Numbered references [1]-[15] are mapped to the keys in ../references.bib
%% (the submitted refs 7 and 9 are two versions of one work and share a key).
%% Build from the parent folder so figures/ and references.bib resolve:
%%   cd .. && latexmk -pdf revision/manuscript_v0_submitted.tex
\documentclass[pdflatex,sn-mathphys-num]{sn-jnl}

\usepackage{graphicx}%
\usepackage{multirow}%
\usepackage{amsmath,amssymb,amsfonts}%
\usepackage{amsthm}%
\usepackage{mathrsfs}%
\usepackage[title]{appendix}%
\usepackage{xcolor}%
\usepackage{textcomp}%
\usepackage{manyfoot}%
\usepackage{booktabs}%
\usepackage{algorithm}%
\usepackage{algorithmicx}%
\usepackage{algpseudocode}%
\usepackage{listings}%

\graphicspath{{figures/}}

\raggedbottom

\begin{document}

\title[CrabNet Hyperparameter Benchmark Dataset]{Materials Science Optimization Benchmark Dataset for High-dimensional, Multi-objective, Multi-fidelity Optimization of CrabNet Hyperparameters}

\author*[1,2]{\fnm{Sterling G.} \sur{Baird}}\email{sterling.baird@byu.edu}

\author[2]{\fnm{Xavier} \sur{Zaitzeff}}

\author[3]{\fnm{Jeet N.} \sur{Parikh}}

\author[1]{\fnm{Taylor D.} \sur{Sparks}}

\affil[1]{\orgdiv{Materials Science \& Engineering},
  \orgname{University of Utah},
  \orgaddress{\street{122 S. Central Campus Drive, \#304},
              \city{Salt Lake City}, \postcode{84112},
              \state{Utah}, \country{USA}}}

\affil[2]{\orgdiv{Mechanical Engineering},
  \orgname{Brigham Young University},
  \orgaddress{\street{701 E University Parkway, 350 EB},
              \city{Provo}, \postcode{84602},
              \state{Utah}, \country{USA}}}

\affil[3]{\orgname{Northwood High School},
  \orgaddress{\street{4515 Portola Pkwy},
              \city{Irvine}, \postcode{92620},
              \state{California}, \country{USA}}}

\abstract{Benchmarks are crucial for driving progress in scientific disciplines. To be effective, benchmarks should closely mimic real-world tasks while being computationally efficient, allowing for accessibility and repeatability. Developing surrogate models indistinguishable from the true observations within the explored dataset bounds dramatically reduces the computational burden of running benchmarks without sacrificing quality, but this requires a large amount of initial data. In material science and chemistry, relevant optimization tasks can be challenging due to their hierarchical, noisy, multi-fidelity, multi-objective, high-dimensional, and non-linearly correlated variables. Additionally, they may include mixed numerical and categorical variables subject to linear and non-linear constraints. Experimentally verifying such tasks can be difficult, making benchmarks essential. This study aimed to overcome these challenges by generating 173219 quasi-random hyperparameter combinations across 23 hyperparameters and using them to train CrabNet on the Matbench experimental band gap dataset (Computational runtime: 387 RTX-2080-Ti GPU days). The results were stored in a free-tier shared MongoDB Atlas dataset, creating a regression dataset that maps hyperparameter combinations to the metrics mean absolute error (MAE), root mean squared error (RMSE), computational runtime, and model size for the trained CrabNet model. To better simulate training, heteroskedastic noise was incorporated into the regression dataset, and bad hyperparameter combinations were excluded. Percentile ranks were computed within each group of identical parameter sets to capture heteroskedastic noise, rather than assuming Gaussian noise as is done in traditional approaches. This approach can be applied to other benchmark datasets, bridging the gap between low cost optimization benchmarks with real-world optimization scenarios.}

\keywords{adaptive design, Bayesian optimization, formulation optimization}

\maketitle

\section*{Specifications table}

\begin{table}[h]
\begin{tabular}{@{}p{0.26\linewidth}p{0.68\linewidth}@{}}
\toprule
\textbf{Subject} & Computational materials science \\
\textbf{Specific subject area} & Composition-based experimental band gap prediction \\
\textbf{Type of data} & Table, Figure \\
\textbf{How the data were acquired} & The data was obtained by executing CrabNet v2.0.8 (available at \url{https://github.com/sparks-baird/CrabNet}) for each of the five folds of the Matbench experimental band gap task (\url{https://matbench.materialsproject.org/Leaderboards\%20Per-Task/matbench_v0.1_matbench_expt_gap/}). Python code in \url{https://github.com/sparks-baird/matsci-opt-benchmarks/blob/7c4346624895a7826ada07ff5e44c2f49eb42b9d/scripts/crabnet_hyperparameter/crabnet_hyperparameter_submitit.py} was used for orchestration, with the University of Utah's Center for High-performance Computing (CHPC) resources. Jobs were sent to the SLURM scheduler using Submitit (\url{https://github.com/facebookincubator/submitit}), and results were logged in JSON format using the MongoDB Data API. The matsci-opt-benchmarks code used for this study can be found at \url{https://github.com/sparks-baird/matsci-opt-benchmarks/tree/v0.2.1} (\url{https://dx.doi.org/10.5281/zenodo.7694289}). \\
\textbf{Data format} & Analyzed, Filtered, Raw \\
\textbf{Description of data collection} & The Ax Platform was used to perform a quasi-random Sobol sampling of 65536 parameter combinations, varying 23 hyperparameters within a constrained search space, with 5 repeats resulting in a total of 327680 training runs. Out of these, 173219 completed successfully, consuming 387 RTX-2080-Ti GPU days or 4614.29 CUDA core years, resulting in 41550 unique sets. To rank repeat simulations, the ``dense'' method with pct=True was used in pandas.core.groupby.GroupBy.rank. \\
\textbf{Data source location} & Free-tier Shared Cluster MongoDB Atlas Database \\
\textbf{Data accessibility} & Repository name: Zenodo. Data identification number: 7694268. Direct URL to data: \url{https://doi.org/10.5281/zenodo.7694268} \\
\bottomrule
\end{tabular}
\end{table}

\section{Objective}

Industry-relevant optimization tasks in the physical sciences are often ``hierarchical, noisy, multi-fidelity~\cite{Existing01_Ghoreishi2019,Existing02_Kandasamy2020_Dragonfly}, multi-objective~\cite{Existing03_Hanaoka2022,Existing04_Hase2018_Chimera}, high-dimensional~\cite{Existing05_Baird2022_CrabNet,Existing06_Eriksson2021_SAASBO}, and non-linearly correlated while exhibiting mixed numerical and categorical variables subject to linear~\cite{Existing07_09_Baird2023_Compactness} and non-linear constraints.''~\cite{Existing07_09_Baird2023_Compactness,Existing08_Baird2023_HardSphere} Existing benchmark datasets~\cite{Existing10_Dunn2020_Matbench,Existing11_DeBreuck2021_MODNet,Existing12_Wang2022_ActiveLearning,Existing13_Liang2021_BO_Benchmarks,Existing14_Henderson2021,Existing15_Hase2021_Olympus}, while very useful, typically are single-objective, single-fidelity, low-dimensional, and ignore or simplify the influence of noise. The aim of this dataset is to better model a physical science optimization task for benchmarking needs while properly quantifying the noise distribution. The inclusion of heteroskedastic noise in a surrogate model avoids assuming a fixed noise distribution and instead build in a more realistic distribution into the surrogate model to virtually match the actual simulation. This approach bridges the gap between low-cost surrogate function evaluations using benchmark datasets and costly, real-world objective function evaluations by considering a multi-objective, multi-fidelity, and high-dimensional task while accounting for heteroskedastic noise.

\section{Data description}

The regression dataset contains hyperparameter sets (including repeats) spanning twenty-three hyperparameters as input for the CrabNet model and their corresponding model's mean absolute error (MAE), root mean squared error (RMSE), computational runtimes, and model size for a fully trained CrabNet model as outputs.

For histogram data for the number of successful repeats see Figure~\ref{fig:repeats}.

For histograms of the mean absolute error, root-mean-square error, runtime, and model size, see Figure~\ref{fig:mae}, Figure~\ref{fig:rmse}, Figure~\ref{fig:gpu}, and Figure~\ref{fig:size}, respectively.

\begin{figure}[h]
\centering
\includegraphics[width=0.6\linewidth]{sobol_repeats_histogram.png}
\caption{Histogram of number of parameter groups vs. number of successful repeats within a given group. The lowest number of repeats for a parameter set is 1, with approximately 2.6 repeats on average.}
\label{fig:repeats}
\end{figure}

\begin{figure}[h]
\centering
\includegraphics[width=0.5\linewidth]{crabnet_mae_histogram.png}
\caption{Histogram of number of training runs vs. mean absolute error using CrabNet on the Matbench experimental band gap task.}
\label{fig:mae}
\end{figure}

\begin{figure}[h]
\centering
\includegraphics[width=0.5\linewidth]{crabnet_rmse_histogram.png}
\caption{Histogram of number of training runs vs. root-mean-square-error using CrabNet on the Matbench experimental band gap task.}
\label{fig:rmse}
\end{figure}

\begin{figure}[h]
\centering
\includegraphics[width=0.5\linewidth]{submitted/crabnet_gpu_runtime_histogram_submitted.png}
\caption{Histogram of number of training runs vs. GPU runtime on an RTX 2080-Ti using CrabNet on the Matbench experimental band gap task. The y-axis is log-scaled.}
\label{fig:gpu}
\end{figure}

\begin{figure}[h]
\centering
\includegraphics[width=0.5\linewidth]{crabnet_model_size_histogram.png}
\caption{Histogram of number of training runs vs. model size using CrabNet on the Matbench experimental band gap task.}
\label{fig:size}
\end{figure}

\section{Experimental design, materials and methods}

A vast number of CrabNet models, totaling to 173219, were trained with different hyperparameter combinations generated by Ax platform's quasi-random Sobol sampling function to ensure uniqueness and a uniform distribution across the search space. While there can be other uses, this dataset is primarily intended as a multi-objective, multi-fidelity, high-dimensional benchmark dataset for formulation-based optimization scenarios by scaling each of the numerical parameters to the range of 0 to 1 and applying a contrived constraint that the sum of all parameters must equal one. Constraining the numerical parameters mimics metal alloy composition constraints found in material science type problems. To realistically capture noise in the benchmark dataset and avoid bias by assuming the noise distribution, simulations were repeated for each quasi-random parameter combination. The repeats were then ranked by the ``dense'' method in pandas.core.groupby.GroupBy.rank and then the rank was added as a feature in the model to predict on for the nondeterministic outputs of MAE, RMSE, and runtime. The additional feature with five repeats then serves as a trained distribution on the experimental data. To improve efficiency and reduce latency, hyperparameter sets (including repeats) were shuffled and divided into batches, then sent to a high-performance computing environment for asynchronous evaluation. Despite some results not being completed due to either timeout or preemption, this trade-off was deemed reasonable for the efficiency and completion gains.

Results were logged to a free-tier MongoDB Atlas database and then aggregated and prepared as machine-learning-ready datasets via Python in Jupyter notebooks. For implementation details, see \url{https://github.com/sparks-baird/matsci-opt-benchmarks/tree/main/scripts/crabnet_hyperparameter} and \url{https://github.com/sparks-baird/matsci-opt-benchmarks/tree/main/notebooks/crabnet_hyperparameter}.

\begin{figure}[h]
\centering
\includegraphics[width=0.49\linewidth]{submitted/shap_mae.png}\hfill
\includegraphics[width=0.49\linewidth]{submitted/shap_rmse.png}\\[4pt]
\includegraphics[width=0.49\linewidth]{submitted/shap_model_size.png}\hfill
\includegraphics[width=0.49\linewidth]{submitted/shap_runtime.png}
\caption{Waterfall SHAP analyses of Histogram Gradient Boosted Regressors trained on each of the different outputs. In clockwise order starting in the top left, MAE, RMSE, runtime, and model size.}
\label{fig:shap}
\end{figure}

\section{Data Sensitivity Analysis}

SHAP analysis as shown in Figure~\ref{fig:shap} analyzes the difference between the true experimental output and the expected output of the model. It measures changes in the expected value given certain variables are fixed experimentally. Looking at the outputs of the model, MAE and RMSE are both model accuracy metrics and we see overlap of impactful hyperparameters between the two. The percentage of data used to train the model decreases predicted MAE and RMSE significantly as well as the emb\_scalar which served as elemental embeddings weights. For model size, the output is a deterministic output based on the number of attention layers, N, and the FC output dimension, d\_model. The model captures this behavior. Considering the runtime output, a larger portion of the data to train on in larger batches, with more parameters to train on over more epochs accounts for most of the difference in overall training runtime.

\section{Value of the data}

\begin{itemize}
\item The dataset is valuable for benchmarking adaptive design methods applied to high-dimensional, constrained, multi-fidelity optimization tasks. Sobol sampling uniformly explores the search space, providing more accurate knowledge of the surface and associated noise through repeats.
\item Practitioners of optimization in the physical sciences can leverage this dataset to simulate real-world materials optimization tasks, such as alloy discovery, and achieve improved performance. This dataset's more real-world characteristics offer better approximations than the commonly used low-dimensional, single objective problems currently used in practice.
\item The dataset is useful in gaining insights into the hyperparameter optimization strategies used for developing compositionally restricted material property prediction models. The normalized and constrained nature of this dataset make it a better approximation for the material science problems which frequently have summation constraints and percentage-based parameters.
\end{itemize}

\section*{Ethics statements}

There are no statements to declare.

\section*{CRediT author statement}

\textbf{Sterling G.~Baird}: Project administration, Conceptualization, Methodology, Software, Validation, Formal analysis, Investigation, Data Curation, Writing - Original Draft, Writing - Review \& Editing, Visualization. \textbf{Jeet N.~Parikh}: Methodology, Software, Formal Analysis, Data Curation, Writing - Original Draft, Writing - Review \& Editing, \textbf{Taylor D.~Sparks}: Supervision, Funding acquisition

\section*{Acknowledgments}

Funding: This work was supported by the National Science Foundation Division of Materials Research [Grant number DMR-1651668].

We thank Trupti Mohanty for discussion related to the future use-case of this dataset as an optimization benchmark. We acknowledge the University of Utah's Center for High Performance Computing (CHPC) for providing computational resources. The code and manuscript associated with a benchmark dataset for particle packing simulations~\cite{Existing08_Baird2023_HardSphere} served as a starting point for the code and manuscript for this work, for which there is shared language and document structure. We acknowledge OpenAI for providing free usage of their research tool, ChatGPT, which was used during the review and editing process.

\section*{Declaration of interests}

The authors declare that they have no known competing financial interests or personal relationships that could have appeared to influence the work reported in this paper.

\bibliography{references}

\end{document}
